% SPDX-License-Identifier: CC-BY-SA-4.0
% Author: Matthieu Perrin

\begingroup

\begin{exercice}[Mini-défi -- Ambiguïté du \og dangling else \fg]\label{exo:parsing/grammars/challenge}

  On considère une grammaire $G$ simplifiée des instructions conditionnelles du C :

  $\begin{array}{rrl}
    \mathit{S} &\rightarrow& \mathit{Instruction} \\
    \mathit{Instruction} & \rightarrow & \text{ ``if(''} \mathit{Condition} \text{ ``)'' } \mathit{Instruction}\\
    &\mid& \text{ ``if(''} \mathit{Condition} \text{ ``)'' } \mathit{Instruction} \text{ ``else'' } \mathit{Instruction} \\
    &\mid& \text{ ``;''} \\
    \mathit{Condition} &\rightarrow& \text{``true''} \mid \text{``false''} \\
  \end{array}$

  \begin{question}
  \item Montrez, en donnant deux arbres de dérivation distincts pour le même mot engendré, que la grammaire $G$ est ambiguë. 
  \item Modifiez la grammaire pour lever l’ambiguïté. 
  \item Justifiez que la nouvelle grammaire n’est pas ambiguë.
  \end{question}

\end{exercice}

\begin{correction}

  \begin{question}
  \item On a par exemple le mot \og{} if(true) if(true) ; else ; \fg{}
    qui peut être interprété comme
    \begin{itemize}
    \item \og{} if(true) \{if(true) ;\} else ; \fg{} ou comme 
    \item \og{} if(true) \{if(true) ; else ;\} \fg.
    \end{itemize}
    
    \begin{tikzpicture}[tree, x=17mm]
      \tree{$S$}{\tree{Instruction}{
          \tree{``if(''}{}
          \tree{Condition}{\tree{``true''}{}}
          \tree{``)''}{}
          \tree{Instruction}{
            \tree[xshift=-1]{``if(''}{}
            \tree{Condition}{\tree[xshift=-1]{``true''}{}}
            \tree[xshift=-1]{``)''}{}
            \tree{Instruction}{\tree[xshift=-1]{``;''}{}}
          }
          \tree[xshift=-2]{``else''}{}
          \tree{Instruction}{\tree[xshift=-2]{``;''}{}}
      }}
    \end{tikzpicture}

    \begin{tikzpicture}[tree, x=17mm]
      \tree{$S$}{\tree{Instruction}{
          \tree{``if(''}{}
          \tree{Condition}{\tree{``true''}{}}
          \tree{``)''}{}
          \tree{Instruction}{
            \tree[xshift=-1]{``if(''}{}
            \tree{Condition}{\tree[xshift=-1]{``true''}{}}
            \tree[xshift=-1]{``)''}{}
            \tree{Instruction}{\tree[xshift=-1]{``;''}{}}
            \tree[xshift=-1]{``else''}{}
            \tree{Instruction}{\tree[xshift=-1]{``;''}{}}
          }
      }}
    \end{tikzpicture}

  \item If faut différencier les instructions associées à un else, et celle qui ne le sont pas :
    $$\begin{array}{rrl}
    \mathit{S} &\rightarrow& A \mid N \\
    \mathit{A} & \rightarrow & \text{ ``if(''} C \text{ ``)'' } A \text{ ``else'' } A  \mid \text{ ``;''} \\
    \mathit{N} & \rightarrow &  \text{ ``if(''} C \text{ ``)'' } N \mid \text{ ``if(''} C \text{ ``)'' } A \text{ ``else'' } N\\
    \mathit{C} &\rightarrow& \text{``true''} \mid \text{``false''} \\
  \end{array}$$

  \item Cette grammaire force systématiquement le rattachement du else au if le plus proche,
    et on peut montrer par induction structurale que toute chaîne dérive soit par $A$
    (zéro else orphelin), soit par $N$ (exactement un chemin avec else orphelin),
    ce qui élimine les dérivations concurrentes.

  \end{question}

\end{correction}

\endgroup
\endinput
